\setcounter{section}{11}






  \zwischenueberschrift{Übungsaufgaben}

\inputaufgabe
{}
{





Es sei $K$ ein
\definitionsverweis {algebraisch abgeschlossener Körper}{}{.}
Beweise
den
Hilbertschen Nullstellensatz

direkt für den
\definitionsverweis {Polynomring}{}{}
in einer Variablen.





}
{} {}

\inputaufgabe
{}
{





Es sei $K$ ein
\definitionsverweis {algebraisch abgeschlossener Körper}{}{}
und

\mavergleichskette
{\vergleichskette
{ f,g
}
{ \in }{ K[X_1 , \ldots , X_n]
}
{  }{
}
{    }{
}
{   }{
}
}
{}{}{.}
Zeige, dass

\mavergleichskette
{\vergleichskette
{V(f)
}
{ \subseteq }{ V(g)
}
{  }{
}
{    }{
}
{   }{
}
}
{}{}{}
genau dann gilt, wenn es eine natürliche Zahl $r$ und ein

\mavergleichskette
{\vergleichskette
{ h
}
{ \in }{ K[X_1 , \ldots , X_n]
}
{  }{
}
{    }{
}
{   }{
}
}
{}{}{}
mit

\mavergleichskette
{\vergleichskette
{ fh
}
{ = }{ g^r
}
{  }{
}
{    }{
}
{   }{
}
}
{}{}{}
gibt. Betrachte auch die Spezialfälle, wo
\mathkor {} {f} {bzw.} {g} {}
konstante Polynome sind.





}
{} {}

\inputaufgabe
{}
{





Zeige, dass sich in der durch
den
Hilbertschen Nullstellensatz

gegebenen Korrespondenz Punkte und
\definitionsverweis {maximale Ideale}{}{}
entsprechen.





}
{} {}

\inputaufgabe
{}
{





Zeige, dass sich in der durch
den
Hilbertschen Nullstellensatz

gegebenen Korrespondenz
\definitionsverweis {irreduzible Varietäten}{}{}
und
\definitionsverweis {Primideale}{}{}
entsprechen.





}
{} {}

\inputaufgabe
{}
{





Es sei $K$ ein
\definitionsverweis {algebraisch abgeschlossener Körper}{}{.}
Beweise den folgenden Spezialfall
des
Hilbertschen Nullstellensatzes

direkt: Wenn

\mavergleichskette
{\vergleichskette
{ f
}
{ \in }{ K[X_1 , \ldots , X_n]
}
{  }{
}
{    }{
}
{   }{
}
}
{}{}{}
keine Nullstelle im $K^n$ besitzt, so ist $f$ ein
\zusatzklammer {von $0$ verschiedenes} {} {}
konstantes Polynom.





}
{} {}

\inputaufgabegibtloesung
{}
{





Wir betrachten die beiden Polynome
\mathl{X^2+Y^2}{} und
\mathl{X^2-Y^3}{} und die zugehörigen algebraischen Kurven über den Körpern $\R$ und ${\mathbb C}$.
\aufzaehlungvierabc{Gilt

\mavergleichskette
{\vergleichskette
{V(X^2+Y^2)
}
{ \subseteq }{ V(X^2-Y^3)
}
{  }{
}
{    }{
}
{   }{
}
}
{}{}{}
in ${\mathbb A}^{2}_{\R}$?
}{Gilt

\mavergleichskette
{\vergleichskette
{V \left(  X^2+Y^2  \right)
}
{ \subseteq }{ V \left(  X^2-Y^3  \right)
}
{  }{
}
{    }{
}
{   }{
}
}
{}{}{}
in ${\mathbb A}^{2}_{ {\mathbb C} }$?
}{Gehört
\mathl{X^2-Y^3}{} zum
\definitionsverweis {Radikal}{}{}
von
\mathl{\left(  X^2+Y^2  \right)}{}  in
\mathl{\R[X,Y]}{?}
}{Gehört
\mathl{X^2-Y^3}{} zum Radikal von
\mathl{\left(  X^2+Y^2  \right)}{}  in
\mathl{{\mathbb C}[X,Y]}{?}
}





}
{} {}

\inputaufgabegibtloesung
{}
{





Es seien
\definitionsverweis {Polynome}{}{}

\mavergleichskette
{\vergleichskette
{ f_1 , \ldots , f_k
}
{ \in }{ {\mathbb C}[X_1 , \ldots , X_n]
}
{  }{
}
{    }{
}
{   }{
}
}
{}{}{}
gegeben, die wir als
\definitionsverweis {Funktionen}{}{}
\maabbdisp {f_i} { {\mathbb C}^n } { {\mathbb C}
} {}
auffassen. Es sei

\mavergleichskette
{\vergleichskette
{ f
}
{ \in }{ {\mathbb C}[X_1 , \ldots , X_n]
}
{  }{
}
{    }{
}
{   }{
}
}
{}{}{}
ein weiteres Polynom und es seien
\maabbdisp {g_1 , \ldots , g_k} { {\mathbb C}^n } { {\mathbb C}
} {}
Funktionen, die nicht unbedingt Polynome sind. Es gelte

\mavergleichskettedisp
{\vergleichskette
{ f
}
{ =} { g_1f_1 + \cdots + g_kf_k
}
{ } {
}
{ } {
}
{ } {
}
}
{}{}{}
\zusatzklammer {eine Gleichung von Funktionen} {} {.}
Zeige, dass $f$ zum
\definitionsverweis {Radikal}{}{}
von
\mathl{\left(  f_1 , \ldots , f_k  \right)}{} gehört.





}
{} {}

\inputaufgabe
{}
{





Es sei $K$ ein
\definitionsverweis {Körper}{}{}
und

\mavergleichskette
{\vergleichskette
{ n
}
{ \in }{ \N_+
}
{  }{
}
{    }{
}
{   }{
}
}
{}{}{.}
Zeige, dass die Funktionen
\maabb {\varphi} { K^n } { K
} {}
der Form

\mavergleichskettedisp
{\vergleichskette
{ \varphi
}
{ =} { { \frac{   P  }{  Q } }
}
{ } {
}
{ } {
}
{ } {
}
}
{}{}{}
mit
\definitionsverweis {Polynomen}{}{}

\mavergleichskette
{\vergleichskette
{ P,Q
}
{ \in }{ K[X_1 , \ldots , X_n]
}
{  }{
}
{    }{
}
{   }{
}
}
{}{}{}
und $Q$ nullstellenfrei auf $K^n$ einen
\definitionsverweis {kommutativen Ring}{}{}
bilden. Zeige, dass bei $K$
\definitionsverweis {algebraisch abgeschlossen}{}{}
dieser mit dem Polynomring übereinstimmt.





}
{} {}

\inputaufgabe
{}
{





Es sei $K$ ein
\definitionsverweis {endlicher Körper}{}{.}
Zeige, dass es nur endlich viele
\definitionsverweis {Nullstellengebilde}{}{}
im
\mathl{{ {\mathbb A}_{  K }^{  n   } }}{} gibt, aber unendlich viele
\definitionsverweis {Radikale}{}{}
in
\mathl{K[X_1 , \ldots , X_n]}{.}





}
{} {}

\inputaufgabe
{}
{





Es sei $R$ ein
\definitionsverweis {kommutativer Ring}{}{}
und sei

\mathbed {f_j} {}
 {j \in J} {}
 {} {} {} {,}
eine Familie von Elementen in $R$. Es sei angenommen, dass die $f_j$ zusammen das
\definitionsverweis {Einheitsideal}{}{}
erzeugen. Zeige, dass es eine endliche Teilfamilie

\mathbed {f_j} {}
 {j \in J_0 \subseteq J} {}
 {} {} {} {,}
gibt, die ebenfalls das Einheitsideal erzeugt.





}
{} {}

\inputaufgabe
{}
{





Es sei $K$ ein
\definitionsverweis {algebraisch abgeschlossener Körper}{}{}
und seien

\mavergleichskette
{\vergleichskette
{ {\mathfrak a} , {\mathfrak b}
}
{ \subseteq }{ K[X_1 , \ldots ,  X_n] }
{  }{}
{    }{}
{   }{}
}
{}{}{}
\definitionsverweis {Radikalideale}{}{.}
Zeige, dass die
\definitionsverweis {Nullstellengebilde}{}{}
$V( {\mathfrak a} )$ und $V( {\mathfrak b} )$ genau dann
\definitionsverweis {affin-linear äquivalent}{}{}
sind, wenn es eine affin-lineare Variablentransformation gibt, die die beiden Ideale ineinander überführt.





}
{} {}







Es sei
\maabbdisp {\varphi} { A } { B
} {}
ein
\definitionsverweis {Ringhomomorphismus}{}{}
zwischen den
\definitionsverweis {kommutativen Ringen}{}{}
\mathkor {} {A} {und} {B} {.}
Zu einem
\definitionsverweis {Ideal}{}{}

\mavergleichskette
{\vergleichskette
{ {\mathfrak a}
}
{ \subseteq }{ A
}
{  }{
}
{    }{
}
{   }{}
}
{}{}{}
nennt man das von
\mathl{\varphi \left(  {\mathfrak a}   \right)}{}
\definitionsverweis {erzeugte Ideal}{}{}
das
\definitionswort {Erweiterungsideal}{}
von ${\mathfrak a}$ unter $\varphi$. Es wird mit
\mathl{{\mathfrak a} B}{} bezeichnet.






Bei $\varphi$ surjektiv ist das einfach das Bildideal.

\inputaufgabe
{}
{





Es sei

\mavergleichskette
{\vergleichskette
{ {\mathfrak a}
}
{ \subseteq }{ \R[X_1 , \ldots , X_n ]
}
{  }{
}
{    }{
}
{   }{
}
}
{}{}{}
ein
\definitionsverweis {Ideal}{}{}
und

\mavergleichskette
{\vergleichskette
{ f
}
{ \in }{ \R[X_1 , \ldots , X_n ]
}
{  }{
}
{    }{
}
{   }{
}
}
{}{}{.}
Zeige, dass genau dann

\mavergleichskette
{\vergleichskette
{ f
}
{ \in }{ {\mathfrak a}
}
{  }{
}
{    }{
}
{   }{
}
}
{}{}{}
gilt, wenn

\mavergleichskette
{\vergleichskette
{ f
}
{ \in }{ {\mathfrak a} {\mathbb C}[X_1 , \ldots , X_n ]
}
{  }{
}
{    }{
}
{   }{
}
}
{}{}{}
für das
\definitionsverweis {Erweiterungsideal}{}{}
gilt.





}
{} {}

\inputaufgabe
{}
{





Skizziere die Graphen der Funktionen $x$ und $y$ auf $V(xy)$.





}
{Man mache sich klar, dass das Produkt $xy$ die Nullfunktion ist.} {}

\inputaufgabe
{}
{





Bestimme den
\definitionsverweis {Koordinatenring}{}{}
zu einer
\definitionsverweis {affin-algebraischen Menge}{}{}

\mavergleichskette
{\vergleichskette
{V
}
{ \subseteq }{ { {\mathbb A}_{  K }^{  n   } }
}
{  }{
}
{    }{
}
{   }{
}
}
{}{}{,}
die aus $d$ Punkten besteht.





}
{} {}

\inputaufgabe
{}
{





Bestimme den
\definitionsverweis {Koordinatenring}{}{}
zur
\definitionsverweis {affin-algebraischen Menge}{}{}

\mavergleichskette
{\vergleichskette
{V
}
{ = }{ V(5X-8Y +3)
}
{ \subseteq }{ {\mathbb A}^{2}_{K}
}
{    }{
}
{   }{
}
}
{}{}{.}





}
{} {}

\inputaufgabe
{}
{





Betrachte die Hyperbel
\mathl{V(xy-1)}{} über dem Körper

\mavergleichskette
{\vergleichskette
{ K
}
{ = }{ \Z/( 11)
}
{  }{
}
{    }{
}
{   }{
}
}
{}{}{.}
Bestimme das Inverse von $4x^3$ im zugehörigen Koordinatenring.





}
{} {}

\inputaufgabe
{}
{





Es sei $K$ ein
\definitionsverweis {Körper}{}{}
und seien

\mavergleichskette
{\vergleichskette
{ V,W
}
{ \subseteq }{ { {\mathbb A}_{  K }^{  n   } }
}
{  }{
}
{    }{
}
{   }{
}
}
{}{}{}
\definitionsverweis {affin-algebra\-ische Mengen}{}{.}
Es sei

\mavergleichskette
{\vergleichskette
{ V
}
{ \subseteq }{ W
}
{  }{
}
{    }{
}
{   }{
}
}
{}{}{}
vorausgesetzt. Man definiere einen
$K$-\definitionsverweis {Algebra\-homomorphismus}{}{}
zwischen den beiden
\definitionsverweis {Koordinatenringen}{}{}
$R(V)$ und $R(W)$ und beschreibe dessen wichtigste Eigenschaften. Man gebe ein Beispiel von zwei affin-algebraischen Mengen, die nicht ineinander enthalten sind, von denen aber die Koordinatenringe isomorph sind.





}
{} {}

\inputaufgabe
{}
{





Es sei $K$ ein
\definitionsverweis {Körper}{}{}
und seien

\mavergleichskettedisp
{\vergleichskette
{ {\mathfrak a} , {\mathfrak b}
}
{ \subseteq} { K[X_1 , \ldots , X_n]
}
{ } {
}
{ } {
}
{ } {
}
}
{}{}{}
\definitionsverweis {Ideale}{}{,}
deren
\definitionsverweis {Radikal}{}{}
übereinstimmt. Zeige, dass es eine natürliche Bijektion zwischen den Radikalen der
\definitionsverweis {Restklassenringe}{}{}

\mathdisp {K[X_1 , \ldots , X_n]/ {\mathfrak a} \text{   und } K[X_1 , \ldots , X_n]/ {\mathfrak b}} {  }

gibt.





}
{} {}

\inputaufgabe
{}
{





Es sei $K$ ein
\definitionsverweis {Körper}{}{}
mit $q$ Elementen und sei

\mavergleichskette
{\vergleichskette
{ V
}
{ = }{ V( {\mathfrak a} )
}
{ \subseteq }{ { {\mathbb A}_{  K }^{  n   } }
}
{    }{
}
{   }{
}
}
{}{}{}
eine
\definitionsverweis {affin-algebraische Menge}{}{.}
Zeige, dass der
\definitionsverweis {Koordinatenring}{}{}
von $V$ nicht gleich
\mathl{K[x_1 , \ldots , x_n]/ \left(  x_1^q-x_1 , \ldots , x_n^q-x_n  \right) + {\mathfrak a}}{} sein muss.





}
{} {}

\inputaufgabe
{}
{





Es sei $K$ ein
\definitionsverweis {Körper}{}{}
der
\definitionsverweis {Charakteristik}{}{}
$0$. Wir betrachten den Schnitt von einem Zylinder und einer Kugel, und zwar

\mavergleichskettedisp
{\vergleichskette
{ C
}
{ =} { V \left(  X^2+Y^2-1  \right)  \cap V \left(  (X-3)^2 +Y^2+Z^2-7  \right)
}
{ \subseteq} { { {\mathbb A}_{  K  }^{  3   } }
}
{ } {
}
{ } {
}
}
{}{}{.}
Zeige, dass man den
\definitionsverweis {Koordinatenring}{}{}
von $C$ als
\definitionsverweis {Restklassenring}{}{}
eines
\definitionsverweis {Polynomrings}{}{}
in zwei Variablen schreiben kann.





}
{} {}





  \zwischenueberschrift{Aufgaben zum Abgeben}

\inputaufgabe
{4}
{





Es sei

\mavergleichskette
{\vergleichskette
{ F
}
{ \in }{ {\mathbb C} [X_1 , \ldots ,  X_n]
}
{  }{
}
{    }{
}
{   }{
}
}
{}{}{}
und sei

\mavergleichskette
{\vergleichskette
{ U
}
{ \subseteq }{ { {\mathbb A}_{  {\mathbb C}  }^{  n   } }
}
{  }{
}
{    }{
}
{   }{
}
}
{}{}{}
eine Teilmenge,  die in der metrischen Topologie offen und nicht leer sei. Es sei

\mavergleichskette
{\vergleichskette
{ F|_{ U}
}
{ = }{ 0
}
{  }{
}
{    }{
}
{   }{
}
}
{}{}{}
die Nullfunktion. Zeige, dass dann $F$ das Nullpolynom ist.





}
{} {}

\inputaufgabe
{3}
{





Beweise

Korollar 11.3

direkt aus

Satz 10.10
.





}
{} {}

\inputaufgabe
{7}
{





Es sei $K$ ein
\definitionsverweis {algebraisch abgeschlossener Körper}{}{}
und $R$ der
\definitionsverweis {Polynomring}{}{}
in $n$ Variablen über $K$. Wir wollen einen alternativen Beweis einsehen, dass

\mavergleichskette
{\vergleichskette
{ \operatorname{Id} \, (V(J))
}
{ = }{ \operatorname{rad} \, (J)
}
{  }{
}
{    }{
}
{   }{
}
}
{}{}{}
für jedes Ideal $J$ in $R$ ist, der auf

Korollar 11.3

aufbaut. Es sei

\mavergleichskette
{\vergleichskette
{ f
}
{ \in }{ \operatorname{Id}\,(V(J))
}
{  }{
}
{    }{
}
{   }{
}
}
{}{}{.}
Betrachte den Ring $R[T]$ und zeige, dass das
\definitionsverweis {Ideal}{}{}

\mavergleichskettedisp
{\vergleichskette
{ J'
}
{ =} { (J, 1 -f \cdot T)
}
{ } {
}
{ } {
}
{ } {
}
}
{}{}{}
trivial ist. Schließe daraus, dass $f$ im Radikal von $J$ liegt.





}
{} {}

\inputaufgabe
{3}
{





Es sei

\mavergleichskette
{\vergleichskette
{ F
}
{ \in }{ K[X_1 , \ldots ,  X_n]
}
{  }{
}
{    }{
}
{   }{
}
}
{}{}{}
und betrachte die dadurch definierte
\definitionsverweis {polynomiale Abbildung}{}{}
\maabbeledisp {\varphi} { { {\mathbb A}_{  K  }^{  n   } } } { { {\mathbb A}_{  K  }^{  n+1   } }
} { \left(  x_1 , \ldots ,  x_n  \right) } { \left(  x_1 , \ldots , x_n, F \left(  x_1 , \ldots , x_n  \right)  \right)
} {,}
die eine Bijektion des
\definitionsverweis {affinen Raumes}{}{}
mit dem
\definitionsverweis {Graphen}{}{}
von $F$ definiert. Zu einer
\definitionsverweis {affin-algebraischen Menge}{}{}

\mavergleichskette
{\vergleichskette
{ V( {\mathfrak a} )
}
{ \subseteq }{ { {\mathbb A}_{  K  }^{  n   } }
}
{  }{
}
{    }{
}
{   }{
}
}
{}{}{}
betrachten wir das Bild

\mavergleichskette
{\vergleichskette
{ V'
}
{ = }{ \varphi(V)
}
{  }{
}
{    }{
}
{   }{
}
}
{}{}{.}
Man zeige, dass $V'$ ebenfalls affin-algebraisch ist und man gebe ein beschreibendes Ideal an. Zeige, dass $V$ genau dann
\definitionsverweis {irreduzibel}{}{}
ist, wenn $V'$ irreduzibel ist.





}
{} {}

\inputaufgabe
{5}
{





Wir betrachten die beiden algebraischen Kurven

\mathdisp {V \left(  x^2+y^2-2  \right) \text{   und } V \left(  x^2+2y^2-1  \right)} {  }

über dem Körper $\Z/( 7 )$. Zeige, dass der Durchschnitt leer ist, und finde einen Erweiterungskörper

\mavergleichskette
{\vergleichskette
{ K
}
{ \supseteq }{ \Z/( 7 )
}
{  }{
}
{    }{
}
{   }{
}
}
{}{}{,}
über dem der Durchschnitt nicht leer ist. Berechne alle Punkte im Durchschnitt über $K$ und über jedem anderen Erweiterungskörper. Man beschreibe auch den Koordinatenring des Durchschnitts.





}
{} {}

\inputaufgabe
{4}
{





Es sei $K$ ein
\definitionsverweis {Körper}{}{}
und seien
\mathl{P_1 , \ldots , P_n}{} endlich viele Punkte in der affinen Ebene ${\mathbb A}^{2}_{ K }$. Es seien

\mavergleichskette
{\vergleichskette
{ a_1 , \ldots , a_n
}
{ \in }{ K
}
{  }{
}
{    }{
}
{   }{
}
}
{}{}{}
beliebig vorgegebene Werte. Zeige, dass es ein Polynom

\mavergleichskette
{\vergleichskette
{ F
}
{ \in }{ K[X,Y]
}
{  }{
}
{    }{
}
{   }{
}
}
{}{}{}
mit

\mavergleichskette
{\vergleichskette
{ F \left(  P_i  \right)
}
{ = }{ a_i
}
{  }{
}
{    }{
}
{   }{
}
}
{}{}{}
für alle

\mavergleichskette
{\vergleichskette
{ i
}
{ = }{ 1 , \ldots , n
}
{  }{
}
{    }{
}
{   }{
}
}
{}{}{}
gibt.





}
{} {}
